% Source: https://artofproblemsolving.com/community/c3h1701961_mathcounts_chapter_2007_target_problem_2
% Stable identity: aops-topic-1701961-q01
After the first fold we have a right isosceles triangle $PQR$ with legs $PQ=PR=8$ and area $\dfrac{1}{2}\cdot8\cdot8=32$.

Folding the right-angle vertex $P$ onto the midpoint $M$ of the hypotenuse $QR$ folds along the perpendicular bisector of $PM$. Since $PM$ is along the axis of symmetry and has length $4\sqrt2$, the fold line passes through the midpoints $X$ and $Y$ of the legs, and the corner triangle $PXY$ (legs $4$, area $8$) flips onto triangle $MXY$ (blue in the figure), which lies inside the remaining trapezoid $XQRY$.

\begin{asy}
usepackage("amsmath");
size(5.5cm);
defaultpen(fontsize(10pt));
pair P=(0,0),Q=(8,0),R=(0,8),X=(4,0),Y=(0,4),Mid=(4,4);
fill(X--Q--R--Y--cycle,gray(0.9));
draw(X--Q--R--Y--cycle,linewidth(0.9));
draw(P--X,dashed);draw(P--Y,dashed);
draw(X--Mid--Y,blue+linewidth(0.8));
dot(P);dot(Mid);
label("$P$",P,SW);label("$Q$",Q,SE);label("$R$",R,NW);
label("$X$",X,S);label("$Y$",Y,W);label("$M$",Mid,NE);
label("$4$",(2,0),S);label("$4$",(6,0),S);label("$4$",(0,2),W);label("$4$",(0,6),W);
label("\mbox{}",(min(currentpicture,user=true)+max(currentpicture,user=true))/2,black); // keeps the text after the figure black when LaTeX typesets the labels inline
\end{asy}

So the figure seen after folding is the trapezoid $XQRY$, with area $32-8=\boxed{24}$ square inches.
